% !TITLE 青玉案·元夕
% !AUTHOR 辛弃疾
% !DYNASTY 两宋
% !GENRE 词
% !ORDER 16

\begin{poem}
东风夜放花千树\textsuperscript{1}，更吹落，星如雨\textsuperscript{2}。
宝马雕车香满路。
凤箫声动\textsuperscript{3}，玉壶光转\textsuperscript{4}，一夜鱼龙舞\textsuperscript{5}。

蛾儿雪柳黄金缕\textsuperscript{6}，笑语盈盈暗香去。
众里寻他千百度\textsuperscript{7}，
蓦然\textsuperscript{8}回首，那人却在，灯火阑珊\textsuperscript{9}处。
\end{poem}

\begin{annotations}
\notetext{1}{花千树：千万棵树开满花——形容元宵节街上张挂的花灯，像一夜春风吹开了千树万树的花朵。}
\notetext{2}{星如雨：烟火像流星雨一样洒落。也可能是形容灯笼密密麻麻如满天星斗。}
\notetext{3}{凤箫声动：排箫的声音响起。凤箫，排箫的美称。}
\notetext{4}{玉壶光转：月亮渐渐西沉（像玉壶在转动）。也有一说指灯笼。}
\notetext{5}{鱼龙舞：鱼形和龙形的灯笼在游动舞蹈。元宵节的灯舞表演。}
\notetext{6}{蛾儿雪柳黄金缕：三种元宵节女子头上的装饰。"蛾儿"是蝴蝶形头饰，"雪柳"是白色柳枝形的绢花，"黄金缕"是金色丝线编成的饰物。都是宋代元宵女子盛装出游的行头。}
\notetext{7}{千百度：千百次、无数次。在人群里找她找了无数次都没找到。}
\notetext{8}{蓦然：忽然、猛然。蓦（mò）。在快要放弃的时候，一回头，看见了她。}
\notetext{9}{阑珊：零落、稀疏。灯火阑珊处——不是最热闹的地方，而是灯火将尽、人渐散去的地方。那个人不在人群中心，而在幽暗冷清处——这是全词的点睛之笔。}
\end{annotations}
